\chapter{Primes in Arithmetic Progressions}
\intro{
    Due to Euler, we can show that there is an infinite number of prime using the fact that the sum of reciprocals of prime $\sum_{p } \frac{1}{p}$ diverge.
A natural refinement is to ask how primes are distributed within specific subsets of integers.
A particularly important case arises when considering arithmetic progressions, sequences of the form:
\begin{equation*}
    a,a+d,a+2d,a+3d,\dots.
\end{equation*}
A natural question is whether such sequences contain infinitely many prime numbers. For example, there are infinitely many primes congruent to $1 \pmod 4$, such as 5,13,17,29,… 5,13,17,29,…, and infinitely many primes congruent to $3 \pmod 4$.

We can try answering this question using previously mentionned methods. Assuming the primes are equally distributed among residu classes generated by coprime integer, we can heuristically conjecture the followings:
\begin{equation*}
    \sum_{\substack{p \le x \\ p \equiv a (q)}}\frac 1p \sim \frac{1}{\phi(q)}\sum_{p \leq x}\frac 1p.
\end{equation*}




This question was answered in 1837 by Dirichlet\parencite{dirichlet_1837}, who proved one of the foundational results of analytic number theory.
\begin{theorem*}[Dirichlet]
    If $(a,d)=1$, then the arithmetic progression $a+nd$ contains infinitely many primes.
\end{theorem*}
}
\section{Dirichlet Character}

\subsection{Definition and basic properties}
\begin{definition}[Dirichlet character]\label{def:dirchar}
    Let $m \ge 1$. We say that $\chi : \mathbb{Z} \to \mathbb{C}$ is a Dirichlet character $\pmod m$ if
    \begin{enumerate}[label={\upshape\roman*)}]
        \item\label{def:dirchar:multiplicity}$\chi(ab) = \chi(a)\chi(b), \text{ } \forall a,b \in \mathbb{Z}$;
        \item\label{def:dirchar:zero_condition}$\chi(a) =0 \Leftrightarrow (a,m) \neq1$;
        \item\label{def:dirchar:periodicity}$\chi(a+m) = \chi(a)$, $\forall a \in \mathbb{Z}$.
    \end{enumerate}
\end{definition}
Using \cref{def:dirchar} \ref{def:dirchar:multiplicity}, we obtain the following: $\chi(1)= \chi(1\cdot1) = \chi(1)\chi(1)$ and $\chi(1)\neq0$, which implies $\chi(1)=1$. By periodicity, $\chi(a)$ only depends on the class of $a \pmod m$, so we can see $\chi$ as a function on $\mathbb{Z}/m\mathbb{Z}$.

\begin{definition}[Principal character]
    We define
    $$\chi_0(n):=
    \begin{cases}
        1 & \text{if } (n,m)=1,\\
        0 & \text{if } (n,m) > 1,
    \end{cases}$$
    the principal character $\pmod m$.
\end{definition}
We usually denote a non-principal character with $\psi$.

It is convenient to work in the more general setting of finite abelian groups. Let $G$ be a finite abelian group. A \emph{character} of $G$ is a group homomorphism $\lambda : G\to\mathbb{C}^\times$. We denote by $\widehat G$ the set of characters of $G$.

\begin{lemma}\label{lem:dirchar_group_char}
    The Dirichlet characters $\pmod m$ are in bijection with the characters of $G = \left(\mathbb{Z}/m\mathbb{Z}\right)^\times$: to $\chi$, we associate its restriction $\lambda(a \bmod m) = \chi(a)$ for $(a,m) = 1$; conversely, to $\lambda\in\widehat G$, we associate $\chi(n) = \lambda(n\bmod m)$ if $(n,m) = 1$ and $\chi(n) = 0$ otherwise.
\end{lemma}
\begin{proof}
    If $\chi$ is a Dirichlet character, its restriction is well defined by periodicity, takes non-zero values by \ref{def:dirchar:zero_condition} and is multiplicative by \ref{def:dirchar:multiplicity}, so it is a homomorphism $G\to\mathbb{C}^\times$. Conversely, let $\lambda \in \widehat G$ and $\chi$ defined as above. Conditions \ref{def:dirchar:zero_condition} and \ref{def:dirchar:periodicity} are clear. For \ref{def:dirchar:multiplicity}: if $(a,m) = (b,m) = 1$, then $(ab,m) = 1$ and $\chi(ab) = \lambda(a)\lambda(b) = \chi(a)\chi(b)$; otherwise $(ab,m)\neq1$ and both sides are $0$. The two maps are inverse of each other since a Dirichlet character is determined by its values on the integers coprime to $m$.
\end{proof}

The key fact on characters is that there are \emph{many} of them.
\begin{lemma}[Extension of characters]\label{lem:extension_char}
    Let $G$ be a finite abelian group and $H$ a subgroup of $G$. Every character of $H$ extends to exactly $[G:H]$ characters of $G$.
\end{lemma}
\begin{proof}
    We proceed by induction on $[G:H]$. If $[G:H] = 1$, there is nothing to prove. Otherwise, let $g\in G\setminus H$, let $H' = \langle H,g\rangle$ and let $n\ge 2$ be the smallest integer such that $g^n\in H$ (it exists since $g^{\lvert G\rvert} = 1 \in H$).

    Every element of $H'$ can be written as $hg^j$ with $h\in H$ and $0\le j<n$, and this writing is unique: if $hg^j = h'g^{j'}$ with $0 \le j' \le j < n$, then $g^{j-j'}\in H$, so $j = j'$ by minimality of $n$, and then $h = h'$. In particular $[H':H] = n$.

    Let $\lambda\in\widehat H$. If $\lambda'$ is an extension of $\lambda$ to $H'$, then $w = \lambda'(g)$ satisfies $w^n = \lambda(g^n)$, and $\lambda'(hg^j) = \lambda(h)w^j$, so $\lambda'$ is determined by $w$. Conversely, let $w$ be one of the $n$ complex roots of $w^n = \lambda(g^n)$ (note $\lambda(g^n)\neq0$) and define $\lambda'(hg^j) = \lambda(h)w^j$, which is well defined by uniqueness of the writing. It is a homomorphism: if $j + j' <n$, it is clear, and if $j + j'\ge n$, then $hg^j\cdot h'g^{j'} = (hh'g^n)g^{j+j'-n}$ and
    \begin{equation*}
        \lambda'\left(hh'g^n g^{j+j'-n}\right) = \lambda(hh')\lambda(g^n)w^{j+j'-n} = \lambda(h)\lambda(h')w^{j+j'} = \lambda'(hg^j)\lambda'(h'g^{j'}).
    \end{equation*}
    So $\lambda$ has exactly $n = [H':H]$ extensions to $H'$. Since $[G:H'] < [G:H]$, by induction each of them has exactly $[G:H']$ extensions to $G$, and two different extensions to $H'$ give different extensions to $G$. Every extension of $\lambda$ to $G$ is obtained this way (by restricting it to $H'$). Thus $\lambda$ has $[H':H][G:H'] = [G:H]$ extensions to $G$.
\end{proof}

\begin{proposition}[Basic properties of Dirichlet character]
    \label{prop:number_dir_char}
    There exist $\phi(m)$ distinct Dirichlet characters $\pmod m$ and the set of all distinct characters $\pmod m$ forms a finite abelian group under multiplication, denoted $\widehat{\left(\mathbb{Z}/m \mathbb{Z}\right)^{\times}}$, whose identity is $\chi_0$. The values $\chi(a)$, for $(a,m) = 1$, are $\phi(m)$-th roots of unity, so $\chi^{-1} = \overline{\chi}$. We also have $\widehat{\left(\mathbb{Z}/m \mathbb{Z}\right)^{\times}} \cong\left(\mathbb{Z}/m \mathbb{Z}\right)^{\times}$.
\end{proposition}
\begin{proof}
    Let $G = \left(\mathbb{Z}/m\mathbb{Z}\right)^\times$, so that $\lvert G\rvert = \phi(m)$. By \cref{lem:dirchar_group_char}, it is enough to study $\widehat G$.

    \emph{Number.} The trivial group $H = \{1\}$ has exactly one character, so by \cref{lem:extension_char}, $\lvert\widehat G\rvert = [G:\{1\}] = \phi(m)$.

    \emph{Group structure.} If $\chi,\chi'\in\widehat G$, then $\chi\chi'(a) := \chi(a)\chi'(a)$ is again a homomorphism, the product is associative and commutative, and $\chi_0$ is the identity. For $(a,m)=1$, Euler's theorem (\cref{thm:euler_theorem}) gives $a^{\phi(m)}\equiv1\pmod m$, so $\chi(a)^{\phi(m)} = \chi(a^{\phi(m)}) = \chi(1) = 1$: $\chi(a)$ is a root of unity, hence $\lvert\chi(a)\rvert = 1$ and $\chi(a)^{-1} = \overline{\chi(a)}$. Thus $\overline\chi$ is a character and $\chi\overline\chi = \chi_0$.

    \emph{Isomorphism.} If $C_n = \langle g\rangle$ is cyclic of order $n$, a character is determined by $\lambda(g)$, which can be any $n$-th root of unity; hence $\lambda \mapsto \lambda(g)$ is an isomorphism $\widehat{C_n}\cong\mu_n\cong C_n$. If $G = G_1\times G_2$, then $(\lambda_1,\lambda_2)\mapsto \left((g_1,g_2)\mapsto\lambda_1(g_1)\lambda_2(g_2)\right)$ is an isomorphism $\widehat{G_1}\times\widehat{G_2}\cong \widehat{G}$ (the inverse is given by restriction to $G_1\times\{1\}$ and $\{1\}\times G_2$). By the structure theorem of finite abelian groups, $G$ is a product of cyclic groups, so $\widehat G\cong G$.
\end{proof}
\noindent\emph{Remark.} The isomorphism $\widehat G\cong G$ is not canonical and will not be used. Only the facts that $\lvert\widehat G\rvert = \lvert G\rvert$ and \cref{lem:existence_dir_char} below will matter.

\subsection{Orthogonality relation}
To prove the orthogonality relation we need the following lemma:
\begin{lemma}\label{lem:existence_dir_char}\leavevmode
    \begin{enumerate}[label={\upshape\roman*)}]
        \item\label{lem:existence_dir_char:i} Let $m\ge 2$ and $a \in \mathbb{Z}$ such that $(a,m)=1$ with $a\not \equiv 1 \pmod m$. There exists $\chi \pmod m$ such that $\chi(a) \neq 1$.
        \item\label{lem:existence_dir_char:ii} For $m\ge 2$, let $\chi$ be a character $\pmod m$ with $\chi \neq \chi_0$. Then there exists $a \in \mathbb{Z}$ with $(a,m) = 1$ such that $\chi(a)\neq 1$ (in particular $\chi(a) \neq 0,1$).
    \end{enumerate}
\end{lemma}
\begin{proof}\leavevmode
    \begin{enumerate}[label={\upshape\roman*)}]
        \item Let $G = \left(\mathbb{Z}/m\mathbb{Z}\right)^\times$ and let $H = \langle a\rangle$ be the subgroup generated by the class of $a$. Since $a\not\equiv1\pmod m$, the order $d$ of $a$ in $G$ satisfies $d\ge2$. The map $\lambda(a^j) = e^{2\pi ij/d}$ is a well defined character of $H$ (because $a^j = a^{j'}$ if and only if $d\mid j-j'$), and $\lambda(a) = e^{2\pi i/d}\neq1$. By \cref{lem:extension_char}, $\lambda$ extends to a character of $G$, which gives, by \cref{lem:dirchar_group_char}, a Dirichlet character $\chi$ with $\chi(a) \neq 1$.
        \item The characters $\chi$ and $\chi_0$ are both $0$ on the integers not coprime to $m$. Since $\chi \neq\chi_0$, they must differ at some $a$ with $(a,m) = 1$, that is $\chi(a)\neq\chi_0(a) = 1$. Also $\chi(a) \neq 0$ by \cref{def:dirchar} \ref{def:dirchar:zero_condition}.
    \end{enumerate}
\end{proof}

\begin{theorem}[Orthogonality relation of Dirichlet character]\leavevmode
    \label{thm:orthogonality_rel}
    \begin{enumerate}[label={\upshape\roman*)}]
        \item $\displaystyle\sum_{a \in \mathbb{Z} / m \mathbb{Z}} \chi(a) = \begin{cases}
                \phi(m)  &\text{ if } \chi= \chi_0,\\
                0 &\text{ if } \chi \neq \chi_0;
            \end{cases}$
        \item\label{item:orthogonality_rel_item:ii} $\displaystyle\sum_{\chi \bmod m} \chi(a) = \begin{cases}
                \phi(m) &\text{ if } a \equiv 1 \pmod m, \\
                0 &\text{ if } a \not\equiv 1 \pmod m.
            \end{cases}$
    \end{enumerate}
\end{theorem}

\begin{proof}\leavevmode
    \begin{enumerate}[label={\roman*)}]
        \item\label{1} Since $\chi(a) = 0$ when $(a,m) \neq 1$, we can restrict the sum to $a \in \left(\mathbb{Z} / m \mathbb{Z}\right)^{\times}$. If $\chi = \chi_0$, every term is $1$ and there are $\phi(m)$ of them:
            \begin{equation*}
                \sum_{a \in \mathbb{Z} / m \mathbb{Z}} \chi_0(a)=  \sum_{a \in \mathbb{Z} / m \mathbb{Z}} \mathds{1}_{(a,m)=1} = \phi(m).
            \end{equation*}
            Now suppose $\chi \neq \chi_0$ and, by \cref{lem:existence_dir_char} \ref{lem:existence_dir_char:ii}, take $\alpha$ with $(\alpha,m) = 1$ and $\chi(\alpha) \neq 1$. We have
            \begin{equation}\label{eq:multiplicity_dir_char}
                \chi(\alpha)\sum_{a \in \left(\mathbb{Z} / m \mathbb{Z}\right)^{\times}} \chi(a) = \sum_{a \in \left(\mathbb{Z} / m \mathbb{Z}\right)^{\times}} \chi(\alpha a).
            \end{equation}
            Since $\alpha$ is invertible $\pmod m$, the map $a \mapsto \alpha a$ is a bijection of $\left(\mathbb{Z} / m \mathbb{Z}\right)^{\times}$ onto itself. This implies:
            \begin{equation}\label{eq:complete_set_residu}
                \sum_{a \in \left(\mathbb{Z} / m \mathbb{Z}\right)^{\times}} \chi(\alpha a)= \sum_{a \in \left(\mathbb{Z} / m \mathbb{Z}\right)^{\times}} \chi(a).
            \end{equation}
            Combining \cref{eq:multiplicity_dir_char} and \cref{eq:complete_set_residu}, we get
            \begin{equation*}
                (\chi(\alpha)-1)\sum_{a \in \left(\mathbb{Z} / m \mathbb{Z}\right)^{\times}} \chi( a) = 0.
            \end{equation*}
            Since we took $\alpha$ such that $\chi(\alpha) \neq 1 $, the sum is $0$.
        \item If $a \equiv 1 \pmod m$, then $\chi(a) = \chi(1) = 1$ for every $\chi$ and, by \cref{prop:number_dir_char},
            \begin{equation*}
                \sum_{\chi \bmod m} \chi(a) = \phi(m).
            \end{equation*}
            If $(a,m)\neq1$, then every term is $0$. Now suppose $(a,m) = 1$ and $a\not \equiv 1 \pmod m$. By \cref{lem:existence_dir_char} \ref{lem:existence_dir_char:i}, there exists $\chi'$ such that $\chi'(a) \neq 1$. We use the same idea as in \ref{1}: since $\chi\mapsto\chi'\chi$ is a bijection of the group of characters onto itself,
            \begin{equation*}
                \chi'(a)\sum_{\chi \bmod m}\chi(a) = \sum_{\chi \bmod m} (\chi' \chi)(a)=\sum_{\chi \bmod m}  \chi(a).
            \end{equation*}
            Thus $(\chi'(a)-1)\sum_{\chi}\chi(a) = 0$ and the sum is $0$.
    \end{enumerate}
\end{proof}
Let $\alpha,\beta\in\mathbb{Z}$ with $(\beta,m) = 1$, and let $\beta^{-1}$ be an inverse of $\beta \pmod m$. Applying \ref{item:orthogonality_rel_item:ii} to $a = \alpha \beta^{-1}$, and using $\chi(\beta^{-1}) = \chi(\beta)^{-1} = \overline\chi(\beta)$, we get an equation that will be more useful later on:
\begin{equation}\label{eq:orthogonality_rel_special_case}
    \sum_{\chi \bmod m} \chi(\alpha ) \overline\chi(\beta) =\begin{cases}
                \phi(m) &\text{ if } \alpha \equiv \beta \pmod m, \\
                0 &\text{ if } \alpha \not\equiv \beta \pmod m.
            \end{cases}
\end{equation}

\section{Dirichlet $L$-functions}
\subsection{Definition and basic properties}
\begin{definition}[$L$-function]\leavevmode\\
    Let $\chi(a)$ a Dirichlet character $\pmod m$, $L(s,\chi):\{s \in \mathbb{C} \mid \Re(s)>1\} \to \mathbb{C}$ we denote the Dirichlet $L$-function:
    \begin{equation*}
        L(s,\chi):= \sum_{n\ge 1}^{}\frac{\chi(n)}{n^s}
    \end{equation*}
\end{definition}
Since $\lvert \chi(a)\rvert\le 1 \text{, }\forall a \in \mathbb{Z}$, we have that the sum is absolutely convergent for $\Re(s) > 1$. Therefore, since $\chi$ is completely multiplicative, can use the Euler product (see \cref{thm:euler-product}) to get:
\begin{equation}\label{eq:L-function_euler_product}
    L(s,\chi)= \prod_{p \text{ prime}}\left(\chi(1) + \frac{\chi(p)}{p^s} + \frac{\chi(p^2)}{p^{2s} }+ \dots\right) = \prod_{p \text{ prime }}\left(1-\frac{\chi(p)}{p^s}\right)^{-1}
\end{equation}
For the principal character, we have the following:
\begin{equation}\label{eq:principal_char_l_function_relation}
    L(s,\chi_0) =   \prod_{p \nmid m}\left(1-\frac{\chi_0(p)}{p^s}\right)^{-1}  = \prod_{p \mid m} \left(1-\frac{1}{p^s}\right) \zeta(s)
\end{equation}
Therefore, $L(s,\chi_0)$ has a simple pole in $s=1$.
\begin{lemma}\label{lem:upper_bound_dir_char_sum}
    Let $\chi \pmod m$ and $A<B \in \mathbb{N}$, then 
    \begin{equation*}
        \left\lvert\sum_{A\le n \le  B}^{}\chi(n)\right\rvert \leq m
    \end{equation*}
\end{lemma}
\begin{proof}
     Let $\alpha = \lfloor\frac{ B-A}{q}\rfloor>0$ . Since $\chi$ is periodic, we have that :
    \begin{equation*}
        \left|\sum_{A\le n \le  B}^{}\chi(n)\right|\leq \left|\sum_{k=0}^{ \alpha-1 } \sum_{j=A +k \alpha}^{A +k \alpha + m} \chi(j)\right|+ \sum_{n=1}^{m} \lvert \chi(n) \rvert.
    \end{equation*}
Using the first orthogonality relation of \cref{thm:orthogonality_rel}, we get that 
\begin{equation*}
    \sum_{j=A +k \alpha}^{A +k \alpha + m} \chi(j) = 0.
\end{equation*}
Thus, 
\begin{equation*}
    \sum_{A\le n \le  B}^{}\chi(n) \leq \sum_{n=1}^{m} \lvert \chi(n) \rvert =m.
\end{equation*}
\end{proof}
\begin{theorem}\label{thm:analytic_continuation_L-function}
    Let $\chi$ be a non-trivial character $\pmod m$, then there exist an analytic continuation of $L(s,\chi),$ $\forall s \in \mathbb{C}$ with $\Re(s)>0$ written:
    \begin{equation}
       L(s,\chi) =  s\int_{1 }^{\infty}\frac{S(t,\chi)}{t^{s+1}}dt,
    \end{equation}
    with $S(t,\chi) = \sum_{n\le t}^{}\chi(t)$.
\end{theorem}
\begin{proof}
    Using Abel Summation formula (see \cref{thm:Abel_Summation_formula}), we get 
    \begin{equation*}
        L(s,\chi)= \lim_{N \to \infty}\sum_{n = 1}^{N}\frac{\chi(n)}{n^s} =\lim_{N \to \infty}\frac{S(N,\chi)}{N^s} -  \int_{1}^{N}\frac{dS(t,\chi)}{t^s}
    \end{equation*}
    By \cref{lem:upper_bound_dir_char_sum}, we know that $\left\lvert S(t,\chi)\right\rvert\le m $, then $\frac{S(N,\chi)}{N^s} \xrightarrow{N \to \infty}0$.
    We then have the following using an integration by part:
    \begin{equation}
        L(s,\chi) = -s \int_{1}^{\infty}\frac{dS(t,\chi)}{t^s} = s \int_{1}^{\infty}\frac{S(t,\chi)}{t^{s+1}}dt
    \end{equation}
    Now we can easily show that the integral is absolutely convergent. Let $s = \sigma + iy$ , we get:
    \begin{equation*}
        \left\lvert L(s,\chi)\right\rvert\leq m \lvert s\rvert \int_{1 }^{\infty }t^{-\sigma - 1}dt = m|s| \left[ \frac{t^{-\sigma}}{-\sigma}\right]^{\infty}_{1}
    \end{equation*} 
    Therefore, this integral is absolutely convergent for $\Re(s) > 0$.
\end{proof}






\begin{theorem}
    \begin{equation*} \label{thm:L-function_neq_0_non-trivial}
        L(1,\psi)\neq 0 \text{, for all non-principal character } \psi \pmod q .
    \end{equation*}
\end{theorem}
\begin{proof}
    Suppose $ L(1,\psi_0)= 0$ for some non-principal character $\psi_0 \pmod q$ and let
\begin{equation*}
    A(s) = \prod_{\chi \pmod q} L(s,\chi)
\end{equation*}
\Cref{eq:principal_char_l_function_relation} tells us that $A(s)$ has a simple pole in $s=1$ that cancels out with the zero of $L(1,\psi_0)$. Therefore, $A(s)$ is analytic for $\Re(s) >0$ by \cref{thm:analytic_continuation_L-function}.
We now want to create a contradiction.We have that
\begin{equation*}
    \log A(s) = \sum_{\chi \pmod q}^{}\log L(s,\chi) .
\end{equation*}
However, by using Euler product (see \cref{thm:euler-product}), and the Taylor expansion of the $\log$, we get that
\begin{equation*}
    \log L(s,\chi)  = \log\prod_{p \nmid q}\left(1-\frac{\chi(p)}{p}\right)^{-1} = -\sum_{p \nmid q}\log\left(1-\frac{\chi(p)}{p}\right)= \sum_{p  \nmid q }^{} \sum_{k \ge 1}^{} \frac{\chi(p)^k}{kp^{ks}}
\end{equation*}
We then use the orthogonality relation (see \cref{thm:orthogonality_rel}) to get the following:
\begin{equation*}
    \log A(s) = \sum_{\substack{p \text{ premier }\\k\ge 1}}^{}\frac{1}{kp^{ks}} \sum_{\chi \tpmod q}^{}\chi(p^k) = \phi(q) \sum_{\substack{p \equiv 1 \tpmod q\\ k\ge 1}}^{} \frac{1}{kp^{ks}}
\end{equation*}
Let $m_p := \ord_q(p)$, then we can write:
\begin{equation*}
    \log A(s) =\phi(q) \sum_{p \nmid q}\sum_{\substack{k \ge 1 \\ m_p \mid k}}^{} \frac{1}{kp^{ks}} = \phi(q)\sum_{p \nmid q}^{} \sum_{j\ge1}^{}\frac{1}{jm_p p^{jm_p s}}
\end{equation*}
Using Taylor expansion of $\log$, we obtain 
\begin{equation*}
    \log A(s) =  \sum_{p \nmid q}^{} \frac{-\phi(q)}{m_p} \log(1 - \frac{1}{p^{m_ps}})
\end{equation*}
Taking the exponential on both sides we get:
\begin{equation}\label{eq:product_l_function}
    A(s) = \prod_{p \nmid q} \left(1- \frac{1}{p^{m_ps}}\right)^{\frac{-\phi(q)}{m_p}}
\end{equation}
Suppose now $\Im(\psi) \neq 0$, then $L(1,\overline{\psi}) = 0$ which implies that $A(s)$
has two zeros in $s=1$ and thus, $A(1)=0$. However, \cref{eq:product_l_function} tells us that $A(1)\ge 1$. This mean that $\exists! \psi \tpmod q$ such that $L(1,\psi)=0$.
\end{proof}



\section{Dirichlet's Theorem on Arithmetic Progressions}
\subsection{Proof of the theorem}
\begin{theorem}[Dirichlet]\label{thm:theorem_dirichlet}
    Let $a,q \in \mathbb{N}$ with $(a,q)=1$. Then, for $x \ge 2$,
    \begin{equation*}
        \sum_{\substack{p\le x \\ p \equiv a (q)}}\frac{1}{p}=  \frac{\log_2 x}{\phi(q)} + O(1),
    \end{equation*}
    where the constant in $O(1)$ depends on $q$.
\end{theorem}
\begin{lemma}\label{lemma:lemma_dirichlet_theorem}
    Let $x\ge2$, $\sigma = 1 + \frac{1}{\log x}$ and $f:\mathbb{N} \to \mathbb{C}$ such that $\lvert f(n)\rvert\le 1$. Then
    \begin{equation*}
        \sum_{p} \frac{f(p)}{p^\sigma} = \sum_{p \le x} \frac{f(p)}{p} + O(1),
    \end{equation*}
    where the constant in $O(1)$ does not depend on $f$ or $x$.
\end{lemma}
\begin{proof}
    We write
    \begin{equation*}
        \sum_{p}\frac{f(p)}{p^\sigma} - \sum_{p\le x}\frac{f(p)}{p} = \sum_{p> x}\frac{f(p)}{p^\sigma} + \sum_{p\le x}f(p)\left(\frac{1}{p^\sigma} - \frac1p\right).
    \end{equation*}
    \emph{The tail.} Since $\lvert f(p)\rvert\le 1$, it is enough to bound $\sum_{p>x}p^{-\sigma}$. We use \cref{thm:Abel_Summation_formula} with $a_n = \mathds{1}_{n \text{ prime}}$, $A = \pi$ and the function $t\mapsto t^{-\sigma}$ on $(x,N]$:
    \begin{equation*}
        \sum_{x<p\le N}\frac{1}{p^\sigma} = \frac{\pi(N)}{N^\sigma} - \frac{\pi(x)}{x^\sigma} + \sigma\int_x^N\frac{\pi(t)}{t^{\sigma+1}}\,dt \le N^{1-\sigma} + \sigma\int_x^\infty\frac{\pi(t)}{t^{\sigma+1}}\,dt.
    \end{equation*}
    Letting $N\to\infty$ and using the Chebyshëv bound $\pi(t)\le K\frac{t}{\log t}$ (\cref{thm:chebyshev}) and $\log t\ge\log x$ for $t \ge x$, we get
    \begin{equation*}
        \sum_{p>x}\frac{1}{p^\sigma}\le K\sigma\int_x^\infty\frac{dt}{t^\sigma\log t}\le\frac{K\sigma}{\log x}\int_x^\infty\frac{dt}{t^\sigma} = \frac{K\sigma}{\log x}\cdot\frac{x^{1-\sigma}}{\sigma-1} = K\sigma e^{-1}\le 2K,
    \end{equation*}
    because $(\sigma-1)\log x = 1$, $x^{1-\sigma} = e^{-1}$ and $\sigma \le 1 + \frac1{\log 2}<3$.

    \emph{The main part.} Using $\lvert f(p)\rvert\le 1$ and the inequality $1-e^{-u} \le u$ for $u \ge 0$, we have
    \begin{equation*}
        \left\lvert \sum_{p \le x} f(p)\left(\frac{1}{p^\sigma} - \frac1p\right) \right\rvert \le\sum_{p \le x} \frac{1-p^{1-\sigma}}{p} = \sum_{ p \le x}\frac{1 - e^{-\frac{\log p}{\log x}}}{p} \le \frac{1}{\log x}\sum_{ p \le x}\frac{\log p}{p }.
    \end{equation*}
    By the first Mertens' estimate (\cref{thm:Mertens} \ref{thm:Mertens:i}), this is $\frac{1}{\log x }\left(\log x + O(1)\right) = O(1)$.
\end{proof}

\begin{proof}[Proof of \Cref{thm:theorem_dirichlet}]
    \emph{Step 1: non-principal characters.} Let $\psi$ be a non-principal character $\pmod q$. We show that
    \begin{equation}\label{eq:sum_non_principal_bound}
        \sum_{p\le x}\frac{\psi(p)}{p} = O(1).
    \end{equation}
    For $\sigma>1$, let $\ell(\sigma,\psi) = \sum_p\sum_{k\ge 1}\frac{\psi(p)^k}{kp^{k\sigma}}$ be the logarithm of $L(\sigma,\psi)$ introduced in the proof of \cref{thm:L-function_neq_0_non-trivial}. The terms with $k\ge2$ are bounded independently of $\sigma$:
    \begin{equation*}
        \left\lvert\sum_p\sum_{k\ge2}\frac{\psi(p)^k}{kp^{k\sigma}}\right\rvert\le\sum_p\sum_{k\ge2}\frac{1}{p^k} = \sum_p\frac{1}{p(p-1)}\le 1,
    \end{equation*}
    so that
    \begin{equation}\label{eq:sum_psi_sigma}
        \sum_p\frac{\psi(p)}{p^\sigma} = \ell(\sigma,\psi) + O(1), \qquad \sigma>1.
    \end{equation}
    We now show that $\ell(\sigma,\psi)$ is bounded for $1<\sigma\le2$. By \cref{thm:analytic_continuation_L-function}, $\sigma\mapsto L(\sigma,\psi)$ is continuous on $[1,2]$. It does not vanish on $(1,2]$ since $L(\sigma,\psi) = e^{\ell(\sigma,\psi)}$, and it does not vanish at $1$ by \cref{thm:L-function_neq_0_non-trivial}. A continuous function $[1,2]\to\mathbb{C}\setminus\{0\}$ has a continuous logarithm $\lambda$ on $[1,2]$, which is bounded since $[1,2]$ is compact. On $(1,2]$, $\ell(\cdot,\psi)$ is continuous (the series converges uniformly on $[1+\delta,2]$ for every $\delta>0$) and $e^{\ell} = e^{\lambda}$, so $\ell-\lambda$ is a continuous function with values in $2\pi i\mathbb{Z}$ on the interval $(1,2]$: it is constant. Hence $\ell(\sigma,\psi)$ is bounded on $(1,2]$.

    Now let $x\ge e$ and $\sigma = 1 + \frac{1}{\log x}\in(1,2]$. By \cref{lemma:lemma_dirichlet_theorem} (with $f = \psi$, since $\lvert\psi(p)\rvert \le 1$) and \cref{eq:sum_psi_sigma},
    \begin{equation*}
        \sum_{p\le x}\frac{\psi(p)}{p} = \sum_p\frac{\psi(p)}{p^\sigma} + O(1) = \ell(\sigma,\psi) + O(1) = O(1).
    \end{equation*}
    For $2\le x<e$, the sum is trivially bounded. This proves \cref{eq:sum_non_principal_bound}.

    \emph{Step 2: detecting the progression.} Since $(a,q)=  1$, the orthogonality relation (\cref{eq:orthogonality_rel_special_case} with $\alpha = p$ and $\beta = a$) gives
    \begin{equation*}
        \sum_{\substack{p \le x \\ p \equiv a (q)}} \frac{1}{p} = \sum_{p\le x}\frac1p\cdot\frac{1}{\phi(q)} \sum_{\chi \bmod q} \overline\chi(a)\chi(p) = \frac{1}{\phi(q)} \sum_{\chi \bmod q}\overline\chi(a) \sum_{p\le x}\frac{\chi(p)}{p}.
    \end{equation*}
    The principal character is like a thorn in our foot, since most of our results don't apply when $\chi = \chi_0$. To manage this situation, we isolate the principal character in our sum:
    \begin{equation*}
        \sum_{\substack{p \le x \\ p \equiv a (q)}} \frac{1}{p} = \frac{1}{\phi(q)} \overline\chi_0(a) \sum_{p\le x}\frac{\chi_0(p)}{p} + \frac{1}{\phi(q)} \sum_{\chi \neq \chi_0}\overline\chi(a) \sum_{p\le x}\frac{\chi(p)}{p}.
    \end{equation*}
    By \cref{eq:sum_non_principal_bound} and $\lvert\overline\chi(a)\rvert = 1$, the second term is a sum of $\phi(q)-1$ bounded terms, so it is $O(1)$. For the first term, $\overline\chi_0(a) = 1$ since $(a,q) = 1$, and $\chi_0(p) = 1$ except for the finitely many primes $p \mid q$. By \cref{thm:Mertens} \ref{thm:Mertens:ii},
    \begin{equation*}
        \sum_{p\le x}\frac{\chi_0(p)}{p} = \sum_{p\le x}\frac1p - \sum_{\substack{p\mid q\\ p \le x}}\frac1p = \log_2 x + O(1).
    \end{equation*}
    We conclude that
    \begin{equation*}
        \sum_{\substack{p \le x \\ p \equiv a (q)}} \frac{1}{p} =\frac{1}{\phi(q)}\log_2 x + O(1).
    \end{equation*}
\end{proof}

\begin{corollary}
    If $(a,d)=1$, then the arithmetic progression $a+nd$ contains infinitely many primes.
\end{corollary}
\begin{proof}
    By \cref{thm:theorem_dirichlet} with $q = d$, $\sum_{p\le x,\ p\equiv a (d)}\frac1p\to\infty$ as $x\to\infty$. If there were only finitely many primes $p \equiv a \pmod d$, this sum would be bounded.
\end{proof}

\section{Prime Number Theorem for Arithmetic Progressions}
\begin{theorem}
    
For $\pi(x,q,a):=\#\{p\le x: p \equiv a \pmod q\}$, 
 \begin{equation*}
    \pi(x,q,a)\sim \frac{1}{\phi(q)}\frac{x}{\log x}, \text{ for }(q,a)=1.
\end{equation*}   
\end{theorem}




\startcontents
